\section{Dirichlet series, Perron inversion, and the contour programme}
\label{sec:contour-mellin}

The paper designated Part I.5 moves the fixed-point numerator
\[
 (1-M_q(n))x-\chi_q(n)
\]
into the coefficients of a Dirichlet series.  Perron inversion can then
recover that numerator by a vertical contour integral.  The coefficient
extraction is sound; the advertised equivalence with every periodic point
inherits the shortest-word defect isolated in
\cref{eq:padded-cycle-code}.

\subsection{The three Dirichlet series}

Let $q$ be an odd prime and $x\in\C$.  With $M_q(n)$ defined using the
shortest binary word of $n\geq1$, put
\begin{align}
 \zeta_{M_q}(s)&=\sum_{n=1}^{\infty}\frac{M_q(n)}{n^s},
 \label{eq:zeta-Mq}\\
 \zeta_{\chi_q}(s)&=\sum_{n=1}^{\infty}\frac{\chi_q(n)}{n^s},
 \label{eq:zeta-chiq}\\
 F_q(s,x)&=\sum_{n=1}^{\infty}
 \frac{(1-M_q(n))x-\chi_q(n)}{n^s}\notag\\
 &=\bigl(\zeta(s)-\zeta_{M_q}(s)\bigr)x-
   \zeta_{\chi_q}(s).
 \label{eq:Fq}
\end{align}
The source writes
\begin{equation}
\label{eq:sigma-q}
 \sigma_q=\log_2\left(\frac{q+1}{2}\right)
\end{equation}
for the common absolute-convergence threshold obtained from its summatory
estimates.  All contour identities below are first asserted on a line
$\Re s=c>\sigma_q$.

\subsection{Dyadic summation}

The digit recursions imply exact block sums.  Write
\[
 A_N=\sum_{0\leq n<2^N}\chi_q(n),
 \qquad
 B_N=\sum_{0\leq n<2^N}M_q(n),
\]
with $M_q(0)=1$.  Separating even and odd indices gives
\begin{align}
 A_{N+1}&=\frac{q+1}{2}A_N+2^{N-1},\label{eq:A-recurrence}\\
 B_{N+1}&=\frac{q+1}{2}B_N+\frac12.
 \label{eq:B-recurrence}
\end{align}
Solving these recurrences yields
\begin{align}
 A_N&=
 \begin{cases}
 N2^N/4,&q=3,\\[2mm]
 \displaystyle\frac{2^N}{q-3}
 \left[\left(\frac{q+1}{4}\right)^N-1\right],&q\ne3,
 \end{cases}
 \label{eq:chi-block-sum}\\
 B_N&=\frac{q\left(\frac{q+1}{2}\right)^N-1}{q-1}.
 \label{eq:M-block-sum}
\end{align}

\begin{proof}
For \cref{eq:A-recurrence}, use
\[
 \chi_q(2m)+\chi_q(2m+1)
 =\frac{q+1}{2}\chi_q(m)+\frac12
\]
and sum over $0\leq m<2^N$.  For \cref{eq:B-recurrence}, the usual
$M_q(2m)=M_q(m)/2$ fails only at $m=0$, because the shortest-word convention
sets $M_q(0)=1$.  Hence the even sum is
$1+\frac12(B_N-1)$ and the odd sum is $qB_N/2$.  The initial values are
$A_0=0$ and $B_0=1$; elementary solution of the recurrences gives the two
closed forms.
\end{proof}

These identities explain the appearance of $\sigma_q$.  A guaranteed
absolute-convergence half-plane for a linear combination is the intersection
of the component half-planes and therefore uses the maximum of their
abscissae, not the source's printed ``minimum.''  The exact abscissa of a
particular linear combination can be smaller when leading terms cancel, so
no equality of abscissae is inferred for arbitrary complex $x$.

\subsection{Perron extraction of one coefficient}

Let
\[
 a_n(x)=(1-M_q(n))x-\chi_q(n).
\]
For $0<a<1$, $n\geq3$, and $c>\sigma_q$, the Perron difference kernel
isolates $a_n(x)$:
\begin{theorem}[coefficient contour identity]
\label{thm:coefficient-contour}
Under the preceding hypotheses,
\begin{equation}
\label{eq:coefficient-contour}
 a_n(x)=\operatorname{PV}\left\{
 \frac1{2\pi i}\int_{c-i\infty}^{c+i\infty}
 \frac{(n+a)^s-(n+a-1)^s}{s}F_q(s,x)\,ds
 \right\}.
\end{equation}
\end{theorem}

\begin{proof}
For a Dirichlet series $F(s)=\sum_{m\geq1}a_mm^{-s}$ absolutely convergent
on $\Re s=c$, Perron's formula gives
\[
 \operatorname{PV}\frac1{2\pi i}
 \int_{c-i\infty}^{c+i\infty}F(s)X^s\frac{ds}{s}
 =\sum_{m<X}a_m
\]
when $X$ is not an integer.  Apply this with $X=n+a$ and with
$X=n+a-1$.  Since $0<a<1$, the two index sets differ only by $m=n$.
Subtracting gives \cref{eq:coefficient-contour}.
\end{proof}

Consequently the integral vanishes exactly when
\begin{equation}
\label{eq:contour-fixed-numerator}
 (1-M_q(n))x=\chi_q(n).
\end{equation}
This is a coefficient theorem.  Its dynamical content is precisely whatever
fixed-point correspondence is valid for the shortest word representing
$n$.

\begin{status}[the printed contour ``weak Collatz'' theorem]
The source states that periodicity is equivalent to the vanishing of
\cref{eq:coefficient-contour} for some $n\geq3$, with separate exclusions:
\begin{enumerate}[label=\textup{(\roman*)}]
\item for $q=3$, $x\in\Z\setminus\{-1,0,1\}$;
\item for $q=5$, $x\in\Z\setminus\{-1,0\}$;
\item for $q\geq7$, apparently $x\in\Z\setminus\{0\}$, although the
      printed set is malformed.
\end{enumerate}
Each part needs its own contour line $c>\sigma_3$, $c>\sigma_5$, or
$c>\sigma_q$; the outer quantifiers in the printed theorem do not state
this correctly.
\end{status}

\begin{sourceissue}[the contour equivalence misses padded cycle words]
The claimed biconditional is not valid for every listed periodic phase.
For example, $2$ is a periodic point of $T_3$ and is not among the source's
exclusions.  If a shortest word fixed $2$, branch correctness would force it
to be a whole number of repetitions of the reversed itinerary $(1,0)$.
Every such word has a high terminal zero and therefore is not shortest.
Equivalently, $2=\chi_3(B_{2,2}(1))$ uses the padded length-two block, while
no shortest positive block yields \cref{eq:contour-fixed-numerator} with
$x=2$.  Thus \cref{thm:coefficient-contour} is canonical, but the full
periodicity equivalence is retained only as a source claim.  A repaired
phase-sensitive contour formulation must index the block length separately
from its ordinary digit sum.
\end{sourceissue}

\subsection{Corrected functional equations for \texorpdfstring{$f_q$ and
$F_q$}{fq and Fq}}

Extend the coefficient sequence by $a_0(x)=0$, and set
\begin{equation}
 \label{eq:Fq-C}
 C_q(x)=\frac{(q-2)x+1}{2}.
\end{equation}
The digit recursions determine both the ordinary and Dirichlet generating
functions without an analytic-continuation premise.

\begin{proposition}[coefficient and ordinary generating-function recursions]
\label{prop:Fq-coefficient-recursion}
For $n\geq1$ and $n\geq0$, respectively,
\begin{equation}
 \label{eq:Fq-coefficient-recursion}
 a_{2n}(x)=\frac{a_n(x)+x}{2},
 \qquad
 a_{2n+1}(x)=\frac q2a_n(x)-C_q(x).
\end{equation}
Consequently, in the formal power-series ring $\C[[z]]$, and analytically
for $|z|<1$,
\begin{equation}
 \label{eq:fq-corrected-functional-equation}
 f_q(z,x):=\sum_{n=1}^{\infty}a_n(x)z^n
 =\frac{qz+1}{2}f_q(z^2,x)
  +\frac{xz^2-((q-2)x+1)z}{2(1-z^2)}.
\end{equation}
\end{proposition}

\begin{proof}
For $n\geq1$, the shortest-word convention and
\cref{eq:chi-even,eq:chi-odd} give
\[
 M_q(2n)=\frac12M_q(n),\qquad
 M_q(2n+1)=\frac q2M_q(n),
\]
and the odd identities remain valid at $n=0$ because $M_q(0)=1$ and
$\chi_q(0)=a_0(x)=0$.  Substitution in
$a_n(x)=(1-M_q(n))x-\chi_q(n)$ proves
\cref{eq:Fq-coefficient-recursion}.  Splitting the formal series into its
positive even indices and its nonnegative odd indices gives
\begin{align*}
 \sum_{n\geq1}a_{2n}(x)z^{2n}
 &=\frac12f_q(z^2,x)+\frac{xz^2}{2(1-z^2)},\\
 \sum_{n\geq0}a_{2n+1}(x)z^{2n+1}
 &=\frac{qz}{2}f_q(z^2,x)-\frac{C_q(x)z}{1-z^2}.
\end{align*}
Their sum is \cref{eq:fq-corrected-functional-equation}.  If the shortest
word of $n$ has length $L$, then
$M_q(n)\leq(q/2)^L$ and
$\chi_q(n)\leq\frac12\sum_{j=0}^{L-1}(q/2)^j$.
Thus $a_n(x)$ has polynomial growth in $n$, so the formal identity
converges absolutely for $|z|<1$.
\end{proof}

\begin{theorem}[corrected Dirichlet functional equation]
\label{thm:Fq-corrected-functional-equation}
Let $q$ be an odd prime, $x\in\C$, and
$\Re s>\sigma_q$.  Then every series below is absolutely convergent and
\begin{equation}
\label{eq:Fq-functional-equation}
\begin{split}
 F_q(s,x)={}&
 \frac{((q-1)x+1)-2^s((q-2)x+1)}
      {2(2^s-2^{\sigma_q})}\,\zeta(s)\\
 &+\frac{q/2}{2^s-2^{\sigma_q}}
 \sum_{m=1}^{\infty}\left(-\frac12\right)^m
 \binom{s+m-1}{m}F_q(s+m,x).
\end{split}
\end{equation}
The series over $m$ converges locally uniformly in that half-plane.
\end{theorem}

\begin{proof}
Put $\alpha=(q+1)/2=2^{\sigma_q}$.  Positivity of $M_q(n)$ and
$\chi_q(n)$, together with \cref{eq:chi-block-sum,eq:M-block-sum}, gives
\[
 \sum_{1\leq n<2^N}|a_n(x)|
 \leq |x|\bigl(2^N+B_N\bigr)+A_N.
\]
For $q>3$ the right side is $O_{q,x}(\alpha^N)$; for $q=3$ it is
$O_x(N2^N)$.  Dyadic subdivision therefore proves absolute convergence of
$F_q(s,x)$ for $\Re s>\sigma_q$.

Split the Dirichlet series by parity.  The even recursion gives
\[
 \sum_{n\geq1}\frac{a_{2n}(x)}{(2n)^s}
 =2^{-s-1}\bigl(F_q(s,x)+x\zeta(s)\bigr).
\]
For the odd part, the binomial series at $1/(2n)$ gives
\begin{equation}
 \label{eq:Fq-odd-shift-expansion}
 \sum_{n\geq1}\frac{a_n(x)}{(2n+1)^s}
 =2^{-s}\sum_{m=0}^{\infty}\left(-\frac12\right)^m
   \binom{s+m-1}{m}F_q(s+m,x).
\end{equation}
To justify the interchange exactly, let $K$ be a compact subset of
$\Re s>\sigma_q$, put
$u=\inf_{s\in K}\Re s$, and choose
$R\geq\max(1,\sup_{s\in K}|s|)$.  From the product formula,
\begin{align*}
 \left|\binom{s+m-1}{m}\right|
 &\leq R\prod_{j=1}^{m-1}\left(1+\frac{R-1}{j+1}\right)\\
 &\leq Re^{R-1}(m+1)^{R-1}.
\end{align*}
Consequently the absolute value of the double series is bounded uniformly
on $K$ by
\[
 Re^{R-1}\left(\sum_{m\geq0}2^{-m}(m+1)^{R-1}\right)
 \left(\sum_{n\geq1}|a_n(x)|n^{-u}\right)<\infty.
\]
This proves both interchanges and local uniformity.  Since
\[
 \sum_{n\geq0}(2n+1)^{-s}=(1-2^{-s})\zeta(s)
\]
the odd recursion in \cref{eq:Fq-coefficient-recursion} and
\cref{eq:Fq-odd-shift-expansion} yield
\begin{align*}
 F_q(s,x)={}&2^{-s-1}\bigl(F_q(s,x)+x\zeta(s)\bigr)\\
 &+\frac q2\,2^{-s}\sum_{m=0}^{\infty}
 \left(-\frac12\right)^m\binom{s+m-1}{m}F_q(s+m,x)\\
 &-C_q(x)(1-2^{-s})\zeta(s).
\end{align*}
Move the $m=0$ and even copies of $F_q(s,x)$ to the left, multiply by
$2^s$, and use
$x/2+C_q(x)=((q-1)x+1)/2$.  This is
\cref{eq:Fq-functional-equation}.
\end{proof}

\begin{sourceissue}[the printed functional equation is false for
\texorpdfstring{$x\ne0$}{x not equal to zero}]
The source substitutes $z/(1-z)$ where the first even correction in
$f_q(z^2,x)$ is $z^2/(1-z^2)$.  Its resulting Dirichlet equation has
zeta coefficient
\[
 -\frac{2^s((q-1)x+1)-1}{2(2^s-2^{\sigma_q})}.
\]
As real $s\to+\infty$, the left side tends to
$a_1(x)=-((q-2)x+1)/2$.  The shifted series is bounded in this limit:
the binomial identity and the estimate
$F_q(s+m,x)=a_1(x)+O(2^{-s-m})$ give
\[
 \sum_{m\geq1}\left(-\frac12\right)^m
 \binom{s+m-1}{m}F_q(s+m,x)
 =a_1(x)\bigl((2/3)^s-1\bigr)+O((2/3)^s).
\]
The printed right side therefore tends to $-((q-1)x+1)/2$.  Thus the
printed identity fails for every $x\ne0$;
\cref{eq:Fq-functional-equation} is the corrected identity.  No
continuation or singularity claim is inherited from the false equation.
A fresh continuation proof from the corrected equation is given below.
The denominator lattice alone still does not establish an actual pole,
because the numerator can vanish.
The coefficient contour identity is unaffected: it uses only the original
absolutely convergent Dirichlet series on $\Re s>\sigma_q$.
\end{sourceissue}

\subsection{Meromorphic continuation from the corrected equation}

Part I.5 uses its shift equation to continue $F_q$, and the public Web II
exposition also reports a meromorphic continuation with a leftward
half-lattice of singularities
\cite{SiegelPartIHalf,SiegelWebII}.  Because the printed equation is false,
the conclusion must be rebuilt from \cref{eq:Fq-functional-equation}.
Put
\begin{align}
 b_m(s)&=\left(-\frac12\right)^m\binom{s+m-1}{m},\label{eq:Fq-bm}\\
 A_q(s,x)&=\frac{((q-1)x+1)-2^s((q-2)x+1)}2,\label{eq:Fq-Aqx}\\
 S_q(s,x)&=\sum_{m=1}^{\infty}b_m(s)F_q(s+m,x).
 \label{eq:Fq-Sqx}
\end{align}

\begin{lemma}[normal convergence of the shifted tail]
\label{lem:Fq-normal-shifted-tail}
For every compact $K\subset\C$, there are constants $C_K,A_K>0$ such
that
\[
 |b_m(s)|\leq C_K2^{-m}(m+1)^{A_K}
 \qquad(s\in K,\ m\geq1).
\]
If $m$ is large enough that $s+m$ lies in the original
absolute-convergence half-plane for every $s\in K$, then
\[
 F_q(s+m,x)=a_1(x)+O_K(2^{-m}).
\]
Hence the tail of \cref{eq:Fq-Sqx} converges normally on $K$.
\end{lemma}

\begin{proof}
Take $R\geq\max(1,\sup_{s\in K}|s|)$.  The product estimate used in the
proof of \cref{thm:Fq-corrected-functional-equation} gives
\[
 \left|\binom{s+m-1}{m}\right|
 \leq Re^{R-1}(m+1)^{R-1},
\]
which proves the first assertion.  Choose $r>\sigma_q$ and put
$D_r=\sum_{n\geq1}|a_n(x)|n^{-r}<\infty$.  If
$u=\inf_{s\in K}\Re s$ and $m+u\geq r$, then
\begin{align*}
 |F_q(s+m,x)-a_1(x)|
 &\leq\sum_{n\geq2}|a_n(x)|n^{-\Re s-m}\\
 &\leq D_r2^{r-u-m}.
\end{align*}
Multiplication by the first bound leaves a convergent exponential--
polynomial majorant, uniformly on $K$.
\end{proof}

\begin{theorem}[global meromorphic continuation]
\label{thm:Fq-global-meromorphic-continuation}
For every odd prime $q$ and every $x\in\C$, the Dirichlet series
$F_q(\,\cdot\,,x)$ has a unique meromorphic continuation to $\C$.
If
\[
 \Lambda_q=\left\{\sigma_q+\frac{2\pi i k}{\log 2}:k\in\Z\right\},
\]
then its poles are contained in the explicitly generated candidate set
\begin{equation}
 \label{eq:Fq-candidate-poles}
 \operatorname{Poles}F_q(\,\cdot\,,x)\subseteq\mathcal P_q,
 \qquad
 \mathcal P_q=
 \{\lambda-n:\lambda\in\Lambda_q,\ n\in\N_0\}
 \cup\{1-n:n\in\N_0\}.
\end{equation}
The inclusion need not be equality: cancellation can make a candidate
singularity removable.
\end{theorem}

\begin{proof}
For $r\in\N_0$, set
\[
 H_r=\{s\in\C:\Re s>\sigma_q-r\}.
\]
The defining series is holomorphic on $H_0$.  Suppose inductively that it
has been continued meromorphically to $H_r$.  If $s\in H_{r+1}$, then
$s+m\in H_r$ for every $m\geq1$.  On a compact subset of $H_{r+1}$,
only finitely many low-shift terms in \cref{eq:Fq-Sqx} can meet poles of
the existing continuation; by
\cref{lem:Fq-normal-shifted-tail}, the remaining holomorphic tail
converges normally.  Thus $S_q(\,\cdot\,,x)$ is meromorphic on
$H_{r+1}$, and
\begin{equation}
 \label{eq:Fq-continuation-recursion}
 F_q(s,x)=\frac{A_q(s,x)\zeta(s)+(q/2)S_q(s,x)}
                  {2^s-2^{\sigma_q}}
\end{equation}
defines a meromorphic extension to $H_{r+1}$.  It agrees with the existing
function on $H_r$, so the identity theorem gives uniqueness on the overlap.
Since $\bigcup_{r\geq0}H_r=\C$, induction proves the continuation.

At the first step, new candidates can occur only at the pole $1$ of
$\zeta$ and at the zero set $\Lambda_q$ of the denominator.  Each later
shifted term replaces an existing candidate $p$ by $p-m$, $m\geq1$,
while division introduces $\Lambda_q$ again.  Induction therefore gives
the containing set in \cref{eq:Fq-candidate-poles}; it does not supply
nonvanishing of the corresponding numerators.
\end{proof}

\subsection{Corrected principal parts for \texorpdfstring{$q=3$}{q=3}}

Write $L=\log 2$ and use $S_3$ as in \cref{eq:Fq-Sqx}.  Then
\begin{equation}
 \label{eq:F3-corrected-continuation}
 F_3(s,x)=\frac{
 \frac12\bigl((2x+1)-2^s(x+1)\bigr)\zeta(s)
 +(3/2)S_3(s,x)}{2^s-2}.
\end{equation}

\begin{proposition}[the corrected $q=3$ principal parts]
\label{prop:F3-corrected-principal-parts}
Let $x\in\C$ and let $\gamma$ be Euler's constant.
\begin{enumerate}[label=\textup{(\roman*)}]
\item At $s=1$, $F_3$ has the genuine double-pole expansion
\begin{equation}
 \label{eq:F3-principal-part-one}
 F_3(s,x)=-\frac1{4L}(s-1)^{-2}
 +\left(
   \frac{3S_3(1,x)}{4L}-\frac x2-\frac38-\frac{\gamma}{4L}
  \right)(s-1)^{-1}+O(1).
\end{equation}
\item If $s_k=1+2\pi ik/L$ with $k\in\Z\setminus\{0\}$, then
$s_k$ is at most a simple pole and
\begin{equation}
 \label{eq:F3-nonreal-lattice-residue}
 \operatorname{Res}_{s=s_k}F_3(s,x)
 =\frac{3S_3(s_k,x)-\zeta(s_k)}{4L}.
\end{equation}
It is removable precisely when $3S_3(s_k,x)=\zeta(s_k)$.
\item The point $s=0$ is a genuine simple pole and
\begin{equation}
 \label{eq:F3-zero-residue}
 \operatorname{Res}_{s=0}F_3(s,x)=-\frac3{16L}.
\end{equation}
\item Every $-k$, $k\in\N_0$, is at most a simple pole.  With
$R_k=\operatorname{Res}_{s=-k}F_3(s,x)$, where $R_k=0$ for a removable
candidate,
\begin{equation}
 \label{eq:F3-negative-residue-recurrence}
 R_k=\frac{3}{16L(1-2^{k+1})(k+1)}
 +\frac{3}{2(1-2^{k+1})}
  \sum_{j=0}^{k-1}2^j\binom{k}{j}R_j.
\end{equation}
The sum is empty for $k=0$.  In particular, every $R_k$ is independent
of $x$; the recurrence alone does not prove $R_k\ne0$ for all $k$.
\end{enumerate}
\end{proposition}

\begin{proof}
Near $s=1$, put $t=s-1$.  Since $S_3$ is holomorphic there,
\begin{align*}
 \frac{(2x+1)-2^s(x+1)}2
   &=-\frac12-(x+1)Lt+O(t^2),\\
 2^s-2&=2Lt+L^2t^2+O(t^3),\\
 \zeta(s)&=t^{-1}+\gamma+O(t).
\end{align*}
Division gives \cref{eq:F3-principal-part-one}; its nonzero
$(s-1)^{-2}$ coefficient proves that the pole is genuinely double.

At $s_k$ with $k\ne0$, both $S_3$ and $\zeta$ are holomorphic,
the denominator has derivative $2L$, and the coefficient of $\zeta$ in
the numerator equals $-1/2$.  This proves
\cref{eq:F3-nonreal-lattice-residue} and the stated removability
criterion.

At $s=0$, only the $m=1$ term of $S_3$ can inherit the double pole at
$1$ to order one:
\[
 -\frac{s}{2}F_3(s+1,x)=\frac1{8L}s^{-1}+O(1).
\]
All $m\geq2$ terms are holomorphic there.  Since $2^0-2=-1$,
\cref{eq:F3-zero-residue} follows.

Finally, fix $k\geq0$ and put $s=-k+t$.  Inductively assume the
candidates $0,-1,\ldots,-(k-1)$ are at most simple.  For
$1\leq m\leq k$,
\[
 \left(-\frac12\right)^m\binom{-k+m-1}{m}
 =2^{-m}\binom{k}{m};
\]
these terms contribute
$\sum_{j=0}^{k-1}2^{-(k-j)}\binom{k}{j}R_j$ to the residue of $S_3$.
For $m=k+1$,
\[
 \left(-\frac12\right)^{k+1}\binom{s+k}{k+1}
 =-\frac{t}{2^{k+1}(k+1)}+O(t^2),
\]
and \cref{eq:F3-principal-part-one} shows that its contribution is
$1/(2^{k+3}(k+1)L)$.  Terms $m\geq k+2$ are holomorphic.  Therefore
\[
 \operatorname{Res}_{s=-k}S_3(s,x)
 =\frac1{2^{k+3}(k+1)L}
  +\sum_{j=0}^{k-1}2^{-(k-j)}\binom{k}{j}R_j.
\]
The zeta term in \cref{eq:F3-corrected-continuation} is holomorphic at
$-k$, and
$2^{-k}-2=(1-2^{k+1})/2^k$.  Multiplication by
$(3/2)/(2^{-k}-2)$ yields
\cref{eq:F3-negative-residue-recurrence} and completes the induction.
\end{proof}

\begin{sourceissue}[the printed $q=3$ pole claims and recurrence]
The source's $q=3$ residues inherit the false $x$-term and therefore
contain spurious factors depending on $4x+1$.  Its headline assertion that
every nonreal lattice point and every nonpositive integer is an actual pole
does not prove noncancellation.  Moreover, the recurrence displayed in the
corollary is undefined at its stated initial index and conflicts with the
formula reached in its proof.  This edition retains the genuine double pole
at $1$, the genuine simple pole at $0$, the exact Laurent coefficients, the
candidate-pole containment, and the corrected recurrence
\cref{eq:F3-negative-residue-recurrence}; it does not turn unproved
noncancellation into a pole classification.
\end{sourceissue}
